%!TEX program = lualatex
% SPDX-License-Identifier: CC-BY-SA-4.0
% Requires: https://github.com/jemmybutton/fiziko (GPL-3.0)
% Set MPINPUTS to the directory containing fiziko.mp before running lualatex.
\documentclass[11pt]{article}

\usepackage{iftex}
\ifLuaTeX\else
  \PackageError{vintage-latex}{This example requires LuaLaTeX}
    {Run lualatex, not pdflatex or xelatex.}
\fi

\usepackage{fontspec}
\usepackage{unicode-math}
\usepackage{xcolor}
\usepackage{geometry}
\usepackage[tracking=true,protrusion=false,expansion=false]{microtype}
\usepackage{luamplib}

\setmainfont{EB Garamond}[Numbers={OldStyle,Proportional}]
\setmathfont{Garamond-Math.otf}
\definecolor{ragpaper}{HTML}{F3E9D6}
\definecolor{printerink}{HTML}{25201B}
\pagecolor{ragpaper}
\color{printerink}
\geometry{a4paper,top=20mm,bottom=20mm,left=24mm,right=24mm}
\pagestyle{empty}
\setlength{\parindent}{1.1em}

\mplibforcehmode
\everymplib{
  input fiziko.mp;
  randomseed := 1969;
  defineMinStrokeWidth(0.2pt);
  color inkcolor; inkcolor := (0.145, 0.125, 0.106);
  drawoptions(withcolor inkcolor);
  beginfig(1);
}
\everyendmplib{endfig;}

\begin{document}

\begin{center}
  {\small\textls[120]{\textsc{Plate III}}}\par
  \vspace{.3em}
  {\Large The Phases of the Moon}\par
  \vspace{.2em}
  {\itshape explained by the position of the moon in its orbit}
\end{center}

\vspace{1.2em}

\begin{center}
\begin{mplibcode}
  % The sun is far to the left. One light direction, kept fixed, shades the
  % earth and every moon alike; the phases follow from geometry alone.
  defineLightDirection(-1/2pi + 0.001, 0);
  numeric orbit, moonDiam;
  orbit := 4.1cm; moonDiam := 0.95cm;

  for i = 0 upto 4:
    drawarrow (-7.3cm, -2.4cm + 1.2cm*i) -- (-5.9cm, -2.4cm + 1.2cm*i)
      withpen pencircle scaled 0.35pt;
  endfor;
  label.top(btex \small\itshape light of the sun etex, (-6.6cm, 2.55cm));

  draw fullcircle scaled 2orbit dashed evenly scaled 0.7
    withpen pencircle scaled 0.3pt;
  draw globe(1.35cm, 40, 20);

  for i = 0 upto 7:
    draw sphere.s(moonDiam) shifted (orbit * dir(180 + 45i));
  endfor;
  label.lft(btex \footnotesize\itshape new etex, (-orbit - 0.6cm, 0));
  label.rt(btex \footnotesize\itshape full etex, (orbit + 0.6cm, 0));
  label.top(btex \footnotesize\itshape first quarter etex, (0, orbit + 0.6cm));
  label.bot(btex \footnotesize\itshape last quarter etex, (0, -orbit - 0.6cm));
\end{mplibcode}
\end{center}

\vspace{.6em}
\begin{center}
  {\small\itshape Fig.\ 1. The moon in eight positions of its orbit, the
  earth in the centre, and the sun's light coming from the left.}
\end{center}

\vspace{1.4em}

\begin{center}
\begin{mplibcode}
  % The same eight moons as they appear from the earth. Now the sphere is
  % fixed before the observer and the light is turned about it instead:
  % fiziko's zero angle lights the face toward the observer (full), and an
  % angle of pi lights the far side (new).
  numeric moonDiam; moonDiam := 1.45cm;
  for i = 0 upto 7:
    defineLightDirection(-pi + 1/4pi*i + 0.001, 0);
    draw sphere.s(moonDiam) shifted (2.15cm*i, 0);
  endfor;
  for i = 0 upto 7:
    draw fullcircle scaled moonDiam shifted (2.15cm*i, 0)
      withpen pencircle scaled 0.3pt;
  endfor;
\end{mplibcode}
\end{center}

\vspace{.4em}
\begin{center}
\small\itshape
\setlength{\tabcolsep}{0pt}
\begin{tabular}{*{8}{p{2.15cm}}}
  \centering new & \centering crescent & \centering first quarter &
  \centering gibbous & \centering full & \centering gibbous &
  \centering last quarter & \centering crescent \tabularnewline
\end{tabular}
\end{center}
\begin{center}
  {\small\itshape Fig.\ 2. The same, as seen from the earth.}
\end{center}

\vspace{1.2em}
\noindent\textls[70]{\textsc{The moon}} has no light of its own, and shows
us only so much of its sunlit half as is turned toward the earth. When it
stands between us and the sun that half is wholly hidden, and the moon is
new; when the earth stands between, the whole of it is seen, and the moon is
full. In the quarters the boundary of light and shade appears as a straight
line, and in the gibbous phases it is convex, though it is at all times the
edge of one and the same circle seen obliquely.

\vfill
\footnotesize\noindent
Both figures use one shading routine. In the first the light is held fixed
and the moons are moved around the earth; in the second the moon is held
before the observer and the light is turned about it. That the two rows agree
is the whole of the explanation, and the agreement is computed rather than
drawn.

\end{document}
